\chapter{Phase 3: Memory and Orchestration (The Brain Phase)}

In this phase, we move beyond stateless interactions to managing context over long-running tasks. We explore how to load just-in-time information (RAG), manage token limits (Compaction), and enforce strict workflows (State Machines).

\section{Step 7 \& 8: RAG and Context Compaction}
We treat the context window like RAM and external storage like a hard drive. Step 7 focuses on retrieving documentation on demand, while Step 8 introduces a method to summarize long histories to stay within token limits.

\begin{lstlisting}[caption={step7\_8\_rag\_compaction.py}, label={lst:step7_8}]
"""
PHASE 3 - STEP 7 & 8: Memory - RAG + Context Compaction
Harness Pillar: Context Management
Goal: Just-in-time retrieval and proactive context management.
"""

import json
import anthropic

KNOWLEDGE_BASE = {"python_async": {"content": "asyncio enables concurrent code..."}}

def simple_retrieval(query: str):
    # Keyword-based retrieval (simplified)
    return [doc for doc in KNOWLEDGE_BASE.values() if any(word in doc['content'] for word in query.split())]

def compact_history(messages: list):
    """Summarize the conversation history down to bullet points."""
    history_text = "\n".join(f"{m['role'].upper()}: {m['content']}" for m in messages[1:-2])
    summary_response = client.messages.create(
        model="claude-opus-4-5-20251101",
        max_tokens=512,
        messages=[{"role": "user", "content": f"Summarize concisely:\n{history_text}"}],
    )
    summary = summary_response.content[0].text
    return [messages[0], {"role": "user", "content": f"[SUMMARY]\n{summary}"}, *messages[-2:]]
\end{lstlisting}

\subsection*{Key Takeaways}
\begin{itemize}
    \item NEVER stuff all documents into the context window upfront; retrieve on demand to save tokens and maintain focus.
    \item Estimate tokens before every API call and compact the history proactively when approaching limits.
    \item Always preserve the original first message and the most recent context during compaction.
\end{itemize}

\section{Step 9: State Machines}
A state machine forces an agent into programmatic, verifiable stages. This prevents the agent from skipping planning, rushing to action, or forgetting to validate its output.

\begin{lstlisting}[caption={step9\_state\_machine.py}, label={lst:step9}]
"""
PHASE 3 - STEP 9: State Machine (Plan -> Execute -> Review)
Harness Pillar: Standard Workflows & Refinement Loops
Goal: Enforce a strict Plan -> Execute -> Review workflow.
"""

from enum import Enum
from dataclasses import dataclass

class State(Enum):
    PLAN = 1; EXECUTE = 2; REVIEW = 3; DONE = 4

@dataclass
class AgentContext:
    task: str; state: State = State.PLAN; plan: list = None

def run_state_machine(task: str):
    ctx = AgentContext(task=task)
    while ctx.state != State.DONE:
        if ctx.state == State.PLAN:
            ctx.plan = ["Step 1", "Step 2"] # Simplified planning
            ctx.state = State.EXECUTE
        elif ctx.state == State.EXECUTE:
            # Execute tools based on plan
            ctx.state = State.REVIEW
        elif ctx.state == State.REVIEW:
            # Separate judge call verifies results
            ctx.state = State.DONE
    return "Task Completed."
\end{lstlisting}

\subsection*{Key Takeaways}
\begin{itemize}
    \item A state machine removes ambiguity from the agentic process by enforcing verified stages.
    \item Use different system prompts for different roles (Planner, Executor, Judge) to maintain professional separation.
    \item The Review stage (LLM-as-judge) should ideally use a separate model call to avoid self-serving bias from the Executor.
\end{itemize}
